The vortex lattice method we shall consider may be interpreted as an inviscid approximation for the case of fully attached, large \textsc{Reynolds} number flows.%
\footnote{%
  \cite{kebbie-athony2019fast} \textsc{Kebbie-Anthony} \emph{et al.} (2019)
}
It has been shown that for at least moderate \textsc{Reynolds} numbers, the fluid flow over a moderately curved airfoil---an object whose treatment is imminent---in two dimensions only depends on the curvature as a first order correction, and that the boundary layer thickness is a function of the \textsc{Reynolds} number.%
\footnote{%
  \cite{lederhilger2024boundary} \textsc{Lederhilger} (2024)
}
We are then led to believe that the thinness of the boundary layer on an airfoil shall be of such parvitude that its direct influence is wholly reduced to the very boundary of the fluid flow.
For a sheet-like thin airfoil, a \emph{lifting surface}, the upper and lower surfaces of the airfoil are one and the same, yielding a sort of double layered boundary layer.
This matter we shall discuss at a later point, but we shall in this section explore manners of approximating such a procedure for any given airfoil.

\subsection{Airfoil theory}
\subfile{airfoil/airfoil.tex}

\subsection{Thin wing approximation}
\subfile{thin/thin.tex}
